\subsection{可测映射与标量可测映射}

若 T 到 Hausdorff 局部凸空间 F 中的映射 f 是标量 \(\mu\)-可测的，一般不能推出 f 是 \(\mu\)-可测的（习题 12）。然而：

命题 12. --- 若 F 是可分的可度量化局部凸空间，则 T 到 F 中的每个标量 \(\mu\)-可测映射 f 都是 \(\mu\)-可测的。

因为， F 可以看作 Banach 空间的可数乘积 \(\prod_{n}E_{n}\) 的子空间（TVS, II, §1, No. 3, 命题 3），并且可以假设 \(\mathrm{pr}_n(\mathbf{F})\) 在 \(\mathbf{E}_n\) 中稠密，后者因此是可分的。对每个 \(n\) ，映射 \(\mathrm{pr}_n \circ \mathbf{f}\) 是标量 \(\mu\)-可测的，因而是 \(\mu\)-可测的（Ch. IV, §5, No. 5, 命题 10 的推论 2），从而 \(\mathbf{f}\) 是 \(\mu\)-可测的（Ch. IV, §5, No. 3, 定理 1）。

命题 13. --- 设 F 是局部凸空间，它是一个由可分、可度量化局部凸空间 \(F_{n}\) 组成的序列的直接极限，且 F 是诸 \(F_{n}\) 之并。令 \(F'\) 为 F 的对偶，赋以拓扑 \(\sigma(\mathbf{F}',\mathbf{F})\) 。则每个标量 \(\mu\)-可测映射 f （ T 到 \(F'\) 中的）都是 \(\mu\)-可测的。

先设 F 是可度量化的且可分的，令 D 为 F 中的可数稠密集。令 \((\mathrm{V}_{n})\) 为 F 中 0 的平衡、凸、开邻域的一个递减基本序列；极集 \(V_{n}^{\circ}\) 是等度连续的，它们的并是整个 \(F'\) 。令 \(\mathrm{T}_{n} = \mathbf{f}^{-1}(\mathrm{V}_{n}^{\circ})\) ；序列 \((\mathrm{T}_{n})\) 是递增的，且 \(T = \bigcup_{n} T_{n}\) ；我们来证明每个 \(T_{n}\) 都是 \(\mu\)-可测的。事实上， \(D \cap V_{n}\) 在 \(V_{n}\) 中稠密；对每个 \(y \in D \cap V_{n}\) ，令 \(S_{y}\) 为 \(t \in T\) 中使 \(|\langle y, f(t) \rangle| \leqslant 1\) 的元素所成之集；假设蕴含每个 \(S_{y}\) 都可测，而 \(T_{n}\) 是诸 \(S_{y}\) \((y \in D \cap V_{n})\) 组成的可数族的交。在此条件下，对 T 的每个紧子集 K 与每个 \(\varepsilon > 0\) ，存在整数 n 使 \(\mu\left(\mathrm{K}-\left(\mathrm{K}\cap\mathrm{T}_{n}\right)\right)\leqslant\frac{\varepsilon}{4}\) ，然后存在紧子集 \(K_{1}\) （含于 \(K\cap T_{n}\) ）使 \(\mu\left(\left(\mathrm{K}\cap\mathrm{T}_{n}\right)-\mathrm{K}_{1}\right)\leqslant\frac{\varepsilon}{4}\) ；最后，存在紧子集 \(K_{2}\) （含于 \(K_{1}\) ）使 \(\mu(K_{1}-K_{2})\leqslant\frac{\varepsilon}{2}\) ，且使在 \(K_{2}\) 上诸函数 \(\langle y,f\rangle\) （ \(y\in D\) ）的限制都连续（Ch. IV, §5, No. 1, 命题 2）。由于集 \(\mathbf{f}(K_{2})\subset\mathbf{f}(T_{n})\subset V_{n}^{\circ}\) 是等度连续的， \(\sigma(F',F)\) 在 \(\mathbf{f}(K_{2})\) 上诱导的拓扑与在 D 中逐点收敛的拓扑相同（GT, X, §2, No. 4, 定理 1）；因此 f 在 \(K_{2}\) 上的限制是连续的，由此得到第一种情形下的断言。

现在过渡到一般情形。若 \(z^{\prime}\) 是 F 上的连续线性形式，其限制 \(z_{n}^{\prime}\) （到 \(F_{n}\) 上）是连续的；由于 \(F = \bigcup_{n} F_{n}\) ，F 的对偶 \(F^{\prime}\) 可以（在代数上）等同于乘积 \(\prod_{n} F_{n}^{\prime}\) 的一个线性子空间，且 \(pr_{n}z^{\prime} = z_{n}^{\prime}\) 。此外，由于 F 的每个有限子集都含于某个 \(F_{n}\) ，拓扑 \(\sigma(F^{\prime}, F)\) 不是别的，正是由诸拓扑 \(\sigma(F_{n}^{\prime}, F_{n})\) 的乘积拓扑所诱导的拓扑。在此条件下，若 f 是标量 \(\mu\)-可测的，则 \(pr_{n} \circ f\) 是标量 \(\mu\)-可测的，因为对每个 \(t \in T\) ， \(\operatorname{pr}_{n}(f(t))\) 是 \(f(t)\) 在 \(F_{n}\) 上的限制。证明的第一部分表明，对每个 n ， \(pr_{n} \circ f\) 是 \(\mu\)-可测的，因此 f 也是（Ch. IV, §5, No. 3, 定理 1）。

